\section{Example 2.1}\label{sec:example2-1}

\subsection{Problem description}\label{subsec:example2-1-problem}

A decision maker chooses among three alternatives A, B and C. The outcome depends on two independent conditions. The value-related condition has two levels \(V_1\) and \(V_2\) with probabilities 0.30 and 0.70. The cost-related condition has two levels \(C_1\) and \(C_2\) with probabilities 0.60 and 0.40. For each alternative, there is a value component associated with the value condition and a cost component associated with the cost condition. Payoff is defined as value minus cost.

The value component for alternatives A, B and C is 7000, 5000 and 4800 under \(V_1\), and 3200, 3600 and 3900 under \(V_2\). The cost component for alternatives A, B and C is 2300, 1600 and 2000 under \(C_1\), and 3500, 2600 and 2700 under \(C_2\). Because the conditions are independent, the combined states of nature are \(S_1\)--\(S_4\) with probabilities 0.18, 0.12, 0.42 and 0.28. Payoffs for each combined state are shown in the payoff table below.

\subsection{Solution}\label{subsec:example2-1-solution}

\subsubsection{(C2) Interpreting the parts of the decision problem}\label{subsubsec:example2-1-c2}

We interpret the narrative and identify the key components.

\begin{center}
	\begin{tabularx}{\linewidth}{@{}L{0.35\linewidth}Y@{}}
		\toprule
		Element & Description \\
		\midrule
		Decision alternatives & Alternative A, Alternative B, Alternative C \\
		States of nature & Combined condition scenarios \(S_1\)--\(S_4\) \\
		Fortuitous events & The actual condition levels that occur \\
		Consequences & Payoff for each alternative and combined scenario \\
		Probabilities & Value condition: 0.30 \(V_1\), 0.70 \(V_2\); Cost condition: 0.60 \(C_1\), 0.40 \(C_2\); independent events \\
		\bottomrule
	\end{tabularx}
\end{center}

\subsubsection{(C3) Building the payoff table from revenue and expense information}\label{subsubsec:example2-1-c3}

From the problem description, we compute the combined-state probabilities and build the payoff table.

We define four combined states of nature:
\(S_1\) is \(V_1\) with \(C_1\), \(S_2\) is \(V_1\) with \(C_2\), \(S_3\) is \(V_2\) with \(C_1\) and \(S_4\) is \(V_2\) with \(C_2\). The joint probabilities are
\[
\begin{aligned}
P(S_1) &= 0.30\cdot 0.60 = 0.18,\\
P(S_2) &= 0.30\cdot 0.40 = 0.12,\\
P(S_3) &= 0.70\cdot 0.60 = 0.42,\\
P(S_4) &= 0.70\cdot 0.40 = 0.28.
\end{aligned}
\]

\begin{center}
	\begin{tabularx}{\linewidth}{@{}L{0.23\linewidth}C C C C@{}}
		\toprule
		State of Nature & Probability & Alternative A & Alternative B & Alternative C \\
		\midrule
		\(S_1\) (\(V_1, C_1\)) & 0.18 & \(7000-2300=4700\) & \(5000-1600=3400\) & \(4800-2000=2800\) \\
		\(S_2\) (\(V_1, C_2\)) & 0.12 & \(7000-3500=3500\) & \(5000-2600=2400\) & \(4800-2700=2100\) \\
		\(S_3\) (\(V_2, C_1\)) & 0.42 & \(3200-2300=900\) & \(3600-1600=2000\) & \(3900-2000=1900\) \\
		\(S_4\) (\(V_2, C_2\)) & 0.28 & \(3200-3500=-300\) & \(3600-2600=1000\) & \(3900-2700=1200\) \\
		\bottomrule
	\end{tabularx}
\end{center}

\subsubsection{(C4) Maximax rule for decision making without probabilities}\label{subsubsec:example2-1-c4}

We ignore probabilities and identify the best payoff each alternative can achieve.

Best payoff for each alternative
\[
\begin{aligned}
\max_i \text{Payoff}(A, S_i) &= 4700\\
\max_i \text{Payoff}(B, S_i) &= 3400\\
\max_i \text{Payoff}(C, S_i) &= 2800
\end{aligned}
\]
\[
\max\{4700, 3400, 2800\} = 4700
\]
The Maximax rule recommends Alternative A.

\subsubsection{(C5) Maximin rule for decision making without probabilities}\label{subsubsec:example2-1-c5}

We focus on the worst payoff for each alternative.

Worst payoff for each alternative
\[
\begin{aligned}
\min_i \text{Payoff}(A, S_i) &= -300\\
\min_i \text{Payoff}(B, S_i) &= 1000\\
\min_i \text{Payoff}(C, S_i) &= 1200
\end{aligned}
\]
We then select the highest of these worst outcomes:
\[
\max\{-300, 1000, 1200\} = 1200
\]
The Maximin rule recommends Alternative C.

\subsubsection{(C6) Maximum opportunity criterion (minimax regret)}\label{subsubsec:example2-1-c6}

We compute regrets by comparing each payoff to the best payoff in each state.

Best payoff in each state
\[
\begin{aligned}
\text{S1} &:\ \max\{4700, 3400, 2800\} = 4700\\
\text{S2} &:\ \max\{3500, 2400, 2100\} = 3500\\
\text{S3} &:\ \max\{900, 2000, 1900\} = 2000\\
\text{S4} &:\ \max\{-300, 1000, 1200\} = 1200
\end{aligned}
\]

Regrets are best payoff minus actual payoff.

\begin{center}
	\begin{tabularx}{\linewidth}{@{}L{0.25\linewidth}C C C@{}}
		\toprule
		State of Nature & Alternative A & Alternative B & Alternative C \\
		\midrule
		\(S_1\) & \(4700-4700=0\) & \(4700-3400=1300\) & \(4700-2800=1900\) \\
		\(S_2\) & \(3500-3500=0\) & \(3500-2400=1100\) & \(3500-2100=1400\) \\
		\(S_3\) & \(2000-900=1100\) & \(2000-2000=0\) & \(2000-1900=100\) \\
		\(S_4\) & \(1200-(-300)=1500\) & \(1200-1000=200\) & \(1200-1200=0\) \\
		\bottomrule
	\end{tabularx}
\end{center}

Maximum regret for each alternative
\[
\begin{aligned}
\max_i \text{Regret}(A, S_i) &= \max\{0, 0, 1100, 1500\} = 1500\\
\max_i \text{Regret}(B, S_i) &= \max\{1300, 1100, 0, 200\} = 1300\\
\max_i \text{Regret}(C, S_i) &= \max\{1900, 1400, 100, 0\} = 1900
\end{aligned}
\]
\[
\min\{1500, 1300, 1900\} = 1300
\]
The minimax regret rule recommends Alternative B.

\subsubsection{(C7) Using probabilities and expected value (Bayes decision rule)}\label{subsubsec:example2-1-c7}

We compute EMV for each alternative by multiplying each payoff by its probability and summing.

For Alternative A
\[
\begin{aligned}
EMV_A &= 0.18(4700) + 0.12(3500) + 0.42(900) + 0.28(-300)\\
&= 846 + 420 + 378 - 84 = 1560
\end{aligned}
\]

For Alternative B
\[
\begin{aligned}
EMV_B &= 0.18(3400) + 0.12(2400) + 0.42(2000) + 0.28(1000)\\
&= 612 + 288 + 840 + 280 = 2020
\end{aligned}
\]

For Alternative C
\[
\begin{aligned}
EMV_C &= 0.18(2800) + 0.12(2100) + 0.42(1900) + 0.28(1200)\\
&= 504 + 252 + 798 + 336 = 1890
\end{aligned}
\]
\[
\max\{1560, 2020, 1890\} = 2020
\]
Bayes EMV recommends Alternative B.

\subsubsection{(C8) Comparing Bayes, Maximin, Maximax and minimax regret}\label{subsubsec:example2-1-c8}

Summary of recommendations

\begin{itemize}
	\item Bayes EMV recommends Alternative B.
	\item Minimax regret recommends Alternative B.
	\item Maximax recommends Alternative A.
	\item Maximin recommends Alternative C.
\end{itemize}

The rules emphasise different aspects of performance: expected value and regret align on B, while Maximax favors the best possible single outcome (A) and Maximin favors safety (C).

\subsubsection{(C9) Tree diagram of the decision-making problem}\label{subsubsec:example2-1-c9}

We represent the combined-state problem as a decision tree.

This diagram summarises the possible combined states and payoffs for each alternative.

\subsubsection{(C10) Evaluating all possibilities in the decision problem}\label{subsubsec:example2-1-c10}

\begin{itemize}
	\item Alternative A is attractive for its best-case payoff but has the poorest worst-case outcome.
	\item Alternative C protects the worst case but has lower expected value than B.
	\item Alternative B is supported by both expected value and minimax regret.
\end{itemize}

Overall, Alternative B offers a balanced decision when probabilities are available.

\subsubsection{(C1) Formulating the final decision and summarising all criteria}\label{subsubsec:example2-1-c1}

Summary across all criteria

\begin{itemize}
	\item Maximax favors Alternative A, prioritising the highest possible payoff.
	\item Maximin favors Alternative C, emphasising protection against the worst case.
	\item Minimax regret favors Alternative B, limiting potential regret.
	\item Bayes EMV favors Alternative B, maximising expected value.
\end{itemize}

Considering all criteria together, Alternative B represents the most coherent final decision because it performs well on expected value and regret and avoids extreme downside risk.

\vspace{6pt}
\noindent\hyperref[table-of-contents]{Back to Top}
